\documentclass[11pt, a4paper]{article}
\usepackage[a4paper, top=2.5cm, bottom=2.5cm, left=2cm, right=2cm]{geometry}
\usepackage{amsmath,amssymb,amsthm,microtype}
\usepackage{hyperref}

\theoremstyle{definition}
\newtheorem{problem}{Problem}
\theoremstyle{plain}
\newtheorem{theorem}{Theorem}
\newtheorem{lemma}{Lemma}

\begin{document}

\title{\vspace{-1.5cm}\textbf{Regional Quantitative Problem-Solving Circle} \\ \Large Problem Set 5: Algorithmic Invariants \& Exchange Arguments (Weeks 9--10)}
\author{\textbf{Lead Architect:} Mohamed Jad Srifi}
\date{\textbf{Term:} Fall 2025 -- Spring 2026}
\maketitle

\noindent\textbf{Instructions:} Do not rely on heuristic guessing. Every problem requires structural justification via formal Exchange Arguments or monotonic potential functions. Ensure all state transitions and partial sums are rigorously evaluated.

\vspace{0.5cm}

\begin{problem}[Interval Scheduling / Adjacent Exchange --- IOI / Kleinberg-Tardos Caliber]
Given a set of $n$ resource requests, where each request $i$ has a start time $s_i$ and a finish time $f_i$, a schedule is valid if no two selected requests overlap in time. 
Prove that the greedy algorithm---which iteratively selects the compatible request with the earliest finish time---is globally optimal and maximizes the total number of scheduled requests.

\textit{(Hint: Define a hypothetical optimal schedule $O$, and execute a mathematical induction over the interval index $r$ to prove that the greedy schedule $G$ finishes no later than $O$ at every step. Conclude by deriving a contradiction on the cardinality $|G| = |O|$.)}
\end{problem}

\vspace{0.5cm}

\begin{problem}[Minimizing Maximum Lateness / Inversion Invariant --- Algorithmic Tripos]
A single processor must execute a set of $n$ independent jobs. Each job $i$ requires a continuous processing time $t_i$ and has a strict deadline $d_i$. The lateness of job $i$ is defined as $L_i = \max(0, C_i - d_i)$, where $C_i$ is its completion time. The goal is to minimize the maximum lateness across all jobs, $L = \max_i L_i$.
Prove that the Earliest Deadline First (EDF) scheduling rule (sorting jobs such that $d_1 \le d_2 \le \dots \le d_n$) is globally optimal.

\textit{(Hint: Assume an optimal schedule contains an inversion---an adjacent pair where $d_i > d_{i+1}$. Prove that swapping this inverted adjacent pair never increases the global maximum lateness.)}
\end{problem}

\vspace{0.5cm}

\begin{problem}[Chip-Firing / Monovariant Termination --- Olympiad Caliber]
Consider a finite line graph with vertices labeled $\{1, 2, \dots, M\}$. Exactly $N$ chips are distributed across these vertices. If any vertex $v$ (where $1 < v < M$) contains at least $2$ chips, a legal "firing" move consists of removing $2$ chips from $v$ and sending exactly one chip to $v - 1$ and exactly one chip to $v + 1$.
Define the potential state function $\Phi = \sum_{v=1}^M v^2 \cdot c(v)$, where $c(v)$ represents the number of chips currently at vertex $v$. 
Using this potential function, prove that the chip-firing game must strictly terminate in a finite number of steps, regardless of the sequence of legal moves chosen.
\end{problem}

\end{document}
